\documentclass[12pt,a4paper]{article}

\usepackage[
  a4paper,
  left=1in,
  right=1in,
  top=1in,
  bottom=1in
]{geometry}
\usepackage{fontspec}
\usepackage{xeCJK}
\usepackage[table]{xcolor}
\usepackage{array}
\usepackage{tabularx}
\usepackage{multirow}
\usepackage{hyperref}
\usepackage{enumitem}
\usepackage{ragged2e}
\usepackage{ulem}
\usepackage{fvextra}

\setmainfont{Segoe UI}[
  Path = C:/Windows/Fonts/,
  Extension = .ttf,
  UprightFont = segoeui,
  BoldFont = segoeuib,
  ItalicFont = segoeuii,
  BoldItalicFont = segoeuiz
]
\setsansfont{Segoe UI}[
  Path = C:/Windows/Fonts/,
  Extension = .ttf,
  UprightFont = segoeui,
  BoldFont = segoeuib,
  ItalicFont = segoeuii,
  BoldItalicFont = segoeuiz
]
\setCJKmainfont{Microsoft YaHei}
\setCJKmonofont{Microsoft YaHei}

\definecolor{wordTitle}{HTML}{000000}
\definecolor{headingGray}{HTML}{404040}
\definecolor{guideBlue}{HTML}{0070C0}
\definecolor{functionalGreen}{HTML}{275317}
\definecolor{functionalBorder}{HTML}{003366}
\definecolor{nonFunctionalRed}{HTML}{660000}

\hypersetup{
  colorlinks=true,
  linkcolor=guideBlue,
  urlcolor=guideBlue
}

\setlength{\parindent}{0pt}
\setlength{\parskip}{0pt}
\renewcommand{\arraystretch}{1.22}
\setlength{\tabcolsep}{0pt}
\pagestyle{empty}

\newcommand{\blankline}{\vspace{\baselineskip}}
\newcommand{\sectitle}[1]{%
  {\fontsize{16pt}{19pt}\selectfont\color{headingGray}#1\par}%
}
\newcommand{\guide}[1]{{\color{guideBlue}#1}}
\newcommand{\tabletext}[1]{\raggedright\arraybackslash #1}
\newcommand{\centerwhite}[1]{\centering\arraybackslash\color{white}#1}
\newcommand{\headercell}[1]{\cellcolor{functionalGreen}\centerwhite{#1}}
\newcommand{\redheadercell}[1]{\cellcolor{nonFunctionalRed}\centerwhite{#1}}
\newcommand{\email}[1]{{\footnotesize #1}}

\begin{document}

{\fontsize{28pt}{33pt}\selectfont IFQ557: Assignment 2 Declaration\par}
\vspace{4pt}

\blankline

\sectitle{Section 1: Team declaration (functional / non-functional)}

\blankline

Please declare your team details on this submission itself by completing and
including one of the two table options below. The details should reflect the
same team as defined in your canvas group. It will assist the marker.

\blankline

\guide{Delete the blue guide text through the document.}

\blankline

\guide{Use either Option A or Option B and delete the one not used.}

\blankline

\textbf{Option A:} \guide{-- Most will use this option. Delete if not used}

\vspace{2pt}

\arrayrulecolor{functionalBorder}
\begin{tabular}{|p{154.5pt}|p{86.0pt}|p{208.25pt}|}
\hline
\multicolumn{3}{|>{\columncolor{functionalGreen}}c|}{\color{white}\textbf{Functional Team -- most teams should use this one}}\\
\hline
\cellcolor{functionalGreen}\centering\color{white}Group/team number:
  & \multicolumn{2}{p{294.25pt}|}{Group 01}\\
\hline
\cellcolor{functionalGreen}\centering\color{white}GitHub or Gitee name and link:
  & \multicolumn{2}{p{294.25pt}|}{\href{https://github.com/Leamon-Lee/IFQ557-Assignment-2}{Leamon-Lee/IFQ557-Assignment-2}}\\
\hline
\headercell{Full name} & \headercell{Student\#} & \headercell{Email}\\
\hline
Rong Yicheng (Bacon) & 2422441046 & \email{rongyicheng23@gmail.com}\\
\hline
Ye Taohao (Leaf) & 2422441048 & \email{2475676977@qq.com}\\
\hline
Li Jingru (Leamon) & 2422441049 & \email{leamon-lee@foxmail.com}\\
\hline
Zhang Zhengyan (Paul) & 2422441050 & \email{zhangzhengyan@pauls-MacBook-Pro.local}\\
\hline
\end{tabular}
\arrayrulecolor{black}

\blankline
\blankline

\newpage

\blankline

\textbf{\color{headingGray}Option B:} {\color{headingGray} }\guide{-- Only use if necessary. Delete if not used}{\color{headingGray}:}\\
To be used when contribution to the assignment was significantly unequal
between team members. Teams are required to collectively negotiate and agree
to a representative \% contribution. All team members must be invited to this
collective discussion and where possible all should agree with the negotiated
outcome (however a majority will be sufficient). This discussion does not have
to be face to face, but this is encouraged. After unanimous (or majority)
agreement has been reached, the team member (nominated by a majority) will
enter the agreed (or majority agreed) percentages into the above table. The
unit coordinator will consider if marks need to be adjusted to reflect
individual contributions to malfunctioning teams. This process is in addition
to the peer assessment process that you will use to rate your teammates.

\blankline
\blankline

\arrayrulecolor{nonFunctionalRed}
\begin{tabular}{|p{145.85pt}|p{83.15pt}|p{167.15pt}|p{51.85pt}|}
\hline
\multicolumn{4}{|>{\columncolor{nonFunctionalRed}}c|}{\color{white}\textbf{\uline{Non-functional} team}}\\
\hline
\cellcolor{nonFunctionalRed}\centering\color{white}Group/team number:
  & \multicolumn{3}{p{304.15pt}|}{Group 01}\\
\hline
\cellcolor{nonFunctionalRed}\centering\color{white}GitHub or Gitee name and link:
  & \multicolumn{3}{p{304.15pt}|}{\href{https://github.com/Leamon-Lee/IFQ557-Assignment-2}{Leamon-Lee/IFQ557-Assignment-2}}\\
\hline
\redheadercell{Full name} & \redheadercell{Student\#} & \redheadercell{Email} & \redheadercell{\%}\\
\hline
Rong Yicheng (Bacon) & 2422441046 & \email{rongyicheng23@gmail.com} & \centering\arraybackslash 28\\
\hline
Ye Taohao (Leaf) & 2422441048 & \email{2475676977@qq.com} & \centering\arraybackslash 47\\
\hline
Li Jingru (Leamon) & 2422441049 & \email{leamon-lee@foxmail.com} & \centering\arraybackslash 22\\
\hline
Zhang Zhengyan (Paul) & 2422441050 & \email{zhangzhengyan@pauls-MacBook-Pro.local} & \centering\arraybackslash 3\\
\hline
 & & & \\
\hline
\multicolumn{3}{|>{\columncolor{nonFunctionalRed}}r|}{\color{white}Total must equal:}
  & \centering\arraybackslash 100\\
\hline
\multicolumn{4}{|p{450.0pt}|}{As the submitter of this assignment, I confirm that the above percentage distribution was agreed by [ all / the majority (correct as appropriate)] of the team members by way of collective negotiation.}\\
\hline
\multicolumn{4}{|p{450.0pt}|}{Additional comments from any team member (list as full name: comment): Percentages in this option are estimated from git commit authorship only and are provided for comparison before choosing the final declaration option.}\\
\hline
\end{tabular}
\arrayrulecolor{black}

\blankline

\guide{In the example above: person 3 contributed less than the others in a 4 person team.}

\blankline
\blankline

\sectitle{Section 2: Team contribution statement}

\blankline

Brief overview from each team member explaining how they contributed: State
your role within the team and the tasks that you completed. There is no word
limit, however, 60-100 words should be sufficient.

\blankline

\begin{tabular}{|p{106.1pt}|p{344.7pt}|}
\hline
Rong Yicheng & Implemented and polished the frontend experience for SoundWave Events, including the homepage prototype, loading entrance page, hero carousel, footer, visual styling, role-aware UI behaviour, organizer dashboard, event display fixes, and ticket page presentation. Also connected frontend expectations to API contracts and helped prepare the interface for final submission.\\
\hline
Ye Taohao & Completed major backend and integration work for authentication, registration, event management, comments, ticket booking, user profile APIs, organizer event status changes, seed data, image handling, password hashing migration, and form submission fixes. Also added and refined tests and corrected several route and pricing bugs during integration.\\
\hline
Li Jingru & Created the initial Flask MVC scaffold, merged the frontend prototype into Flask views/templates, wrote README and submission documentation, developed the value object layer and validation rules, refactored models to use value objects, connected frontend API routes, and stabilized database rebuild/startup behaviour.\\
\hline
Zhang Zhengyan & Built the early database and model foundation, including core user, participant, and organizer model definitions supporting login and registration tests. This work helped establish the SQLAlchemy data layer that later features, services, and tests extended.\\
\hline
\end{tabular}

\blankline

\sectitle{Section 3: git shortlog > output.txt}

\blankline

\begin{Verbatim}[fontsize=\scriptsize,breaklines=true]
Leamon Lee (16):
      Initial Flask MVC scaffold
      Merge bacon frontend into Flask views
      Add project README and UML guide
      Rename site to SoundWave Events
      Refactor Paul models with value objects
      Apply value objects across domain models
      Require value objects in user roles
      Complete value object constraints from UML
      Add detailed submission notes
      Merge leamon value objects into temp
      Refactor models to use value objects
      index on main: 72a6bad 通过pytest进行的项目测试，测试结果报告位于test文件夹中
      On main: preserve local deletions before switching to bacon
      Add database rebuild toggle to runner
      Connect frontend API routes
      Stabilize database rebuild runner

rongyicheng23-afk (20):
      Add SoundWave event prototype
      Refine homepage hero carousel
      Add polished site footer
      Prepare frontend for final submission
      Merge remote-tracking branch 'origin/main' into bacon
      Add SoundWave loading entrance page
      Refine loading page with dynamic concert hero
      Polish loading page typography
      Increase loading page background impact
      Use custom generated loading background
      Replace loading video with smooth realtime ambience
      Refine loading page ambience without line effects
      Add dramatic animated concert light show
      Optimize loading light show performance
      Refine hero actions and auto-dismiss alerts
      Add role-based auth UI and API contract
      Make loading page default entry
      Refine organizer dashboard UI
      Fix event page display issues
      Match ticket pages to event visuals

yetaohao (30):
      登录，注册，登出功能完善，修改了login,signup部分前端代码
      前端页面修改:原数据写死改为从数据库中获取动态数据,新增错误处理页面
      根据最新models完善用户认证和事件管理功能，重构相关模型和路由，更新模板以提升用户体验
      对上一步提交的一些补充
      Revert "对上一步提交的一些补充"
      Revert "根据最新models完善用户认证和事件管理功能，重构相关模型和路由，更新模板以提升用户体验"
      根据Paul最新版models完善用户模型和参与者模型，添加事件创建、更新、取消及注册功能；更新认证路由以支持用户登录和注册
      修改了一个冲突内容
      错误拉取主分支后出现报错，修正提交
      feat: 新增评论系统和购票预订流程
      feat: 新增创建/编辑活动功能
      feat: 填充全部空壳页面
      feat: 实现首页搜索/筛选和Organizer公告功能
      feat: 新增Organizer取消活动功能
      修正了文本模板切片操作失败的问题
      修复图片缺失问题
      通过pytest进行的项目测试，测试结果报告位于test文件夹中
      清理html文件垃圾，修改“启动重建数据库”为关
      统一密码哈希方案（werkzeug → flask_bcrypt),同时修改了上一次在main提交中错误提交的一个路由跳转bug
      新增 Admin API 端点
      新增登出 API
      实现模型层空壳方法
      Web 表单登录路由增加角色验证,Web 表单注册路由支持创建 Organizer
      新增用户资料管理 API
      新增组织者修改活动状态
      修改种子数据，新增更多活动
      活动图片系统改造
      修正了组织者创建活动提交表单错误问题，补全创建活动决定票价的功能
      在 form = EditEventForm(obj=event) 之后加一行 form.ticket_price.data = event._ticket_price，直接将底层的 Decimal 值赋给表单字段。POST 提交时 validate_on_submit() 会用提交数据覆盖该值，不影响写入逻辑。
      修复了修改价格不能够正确返回的bug

yetaohao4-cmd (4):
      Merge branch 'main' into Leaf
      Merge pull request #4 from Leamon-Lee/Leaf
      Merge pull request #5 from Leamon-Lee/Leaf
      Merge pull request #6 from Leamon-Lee/Leaf

张峥岩 (2):
      feat: 完成 User, Participant, Organizer 模型定义，支持登录注册测试
      feat: 完成所有数据库
\end{Verbatim}

\blankline

\sectitle{Section 4: Team use of AI statement}

\blankline

Instructions:
\begin{itemize}[leftmargin=.5in, label=\textbullet, topsep=0pt, itemsep=8pt, parsep=0pt]
  \item We did not use AI. Include part A. Delete part B.
  \item We (one or more team members) used AI in an authorised way (as described in the assignment brief). Complete part B. Delete part A.
\end{itemize}

\blankline

Part A : (delete part if not used)

\blankline

Delete this part if Part B is used.

\blankline
\blankline

Part B : (delete part if not used)

\blankline

As a team, we report that some members used AI in the following authorized way.

\blankline

Summary of AI use

\blankline

AI tools were used as authorised development support for brainstorming,
debugging, refactoring, documentation drafting, LaTeX formatting, and test
interpretation. Team members used AI suggestions to improve wording, identify
possible implementation issues, and speed up repetitive coding tasks, but the
final code, project structure, tests, and submission content were reviewed,
adapted, and accepted by the team members before submission.

\blankline
\blankline

Key examples of AI use

Your marker needs to clearly understand how your team used AI to improve your
work. Provide representative samples of how you used AI using the requested
format.

(Repeat for each example. There is no word limit for this declaration. Please
also include the separator line)

\{
\begin{itemize}[leftmargin=.5in, label=\textbullet, topsep=0pt, itemsep=8pt, parsep=0pt]
  \item Prompt: ``Based on this Flask music event project, help summarize each team member's contribution from the git history and repository files.''
  \item Input/upload: Repository files including README, SUBMIT notes, tests, routes, models, templates, and git shortlog output.
  \item Output used: Draft contribution summaries for Bacon, Leaf, Leamon, and Paul, which were checked against the repository and edited into this declaration.
\end{itemize}

{\footnotesize ========================= NEXT SAMPLE OF AI USE ====================}

\}

\end{document}
